\DocumentMetadata{lang=es-DO,testphase={phase-III,math}}
% Plantilla de PRUEBÍN -- evaluación corta calificada (20-30 min, una o dos
% preguntas). Ver STYLE.md § Reglas por tipo de documento.
%
% Sobre scrartcl, NO sobre la clase exam: un pruebín no necesita puntaje
% automático ni selección múltiple, y sobre scrartcl sí se conserva el PDF
% etiquetado (\DocumentMetadata), que exam.cls obliga a sacrificar.
%
% Diferencia con `assignment`: el pruebín se resuelve en el aula, así que
% lleva cabecera de fecha/duración/puntaje y línea de Nombre/ID, y NO lleva
% bloque de objetivos de aprendizaje.
\documentclass[11pt]{scrartcl}
\usepackage{icv}

% Clave de respuestas. Dos formas de activarla:
%   1. Descomentar \toggletrue{icvclave} aquí abajo, o
%   2. Sin tocar el archivo, desde la raíz del repo:
%      latexmk -pretex='\def\icvclaveon{}' -usepretex \
%              -jobname=<nombre>-clave templates/pruebin/pruebin.tex
\newtoggle{icvclave}
\ifdefined\icvclaveon\toggletrue{icvclave}\fi
%\toggletrue{icvclave}

% -----------------------------------------------------------------------
% Metadata -- reemplaza estos valores para cada pruebín nuevo.
% -----------------------------------------------------------------------
\icvsetup{
  asignatura  = {ICV 441-T -- Hidráulica Aplicada},
  titulocorto = {Pruebín N -- Tema},
  titulo      = {Pruebín N},
  subtitulo   = {Tema específico que se evalúa (opcional)},
  profesor    = {Dr. Ricardo Hernández Moreira},
  periodo     = \icvciclo{2026}{3},
  version     = {v1},
}

\begin{document}
\icvmaketitle

% -- Cabecera del pruebín: fecha, duración y puntaje en una sola línea. --
\vspace{-6pt}
\begin{center}\small
  \textbf{Fecha:} TODO: falta fecha \;\textbullet\;
  \textbf{Duración:} TODO: falta duración \;\textbullet\;
  \textbf{Puntaje total:} TODO: falta puntaje total
\end{center}

\vspace{-4pt}
\makebox[\linewidth][l]{%
  \textbf{Nombre:}\hspace{0.5em}\rule{0.46\linewidth}{0.4pt}%
  \hspace{1em}\textbf{ID:}\hspace{0.5em}\hrulefill%
}

\begin{icvnote}[Instrucciones]
Material permitido, valor de las constantes que deben usar, formato de
respuesta y cualquier otra regla. Si el pruebín pide un trazado o un
dibujo, especifica aquí cómo distinguir cada línea \textbf{por trazo y por
rótulo de texto}, nunca solo por color --- ver STYLE.md § Accesibilidad
para daltonismo.
\end{icvnote}

\begin{icvproblem}[points=20]{Título del problema}
Enunciado. Los datos numéricos van en una tabla con columnas \texttt{S} de
siunitx; los números, siempre con \verb|\num|/\verb|\qty|, nunca con formato
manual.

% Nota: bajo PDF etiquetado (phase-III), el entorno enumerate NO acepta las
% claves de enumitem (label=, itemsep=...). Se usan etiquetas explícitas.
\begin{enumerate}
  \item[\textbf{(a)}] \textbf{(8 puntos)} Primer inciso.
  \item[\textbf{(b)}] \textbf{(4 puntos)} Segundo inciso.
  \item[\textbf{(c)}] \textbf{(8 puntos)} Tercer inciso.
\end{enumerate}
\end{icvproblem}

% -- Anexo gráfico opcional (diagrama, retícula de trazado, tabla para
%    llenar a mano). Va en su propia página; requiere \usepackage{tikz} en
%    el preámbulo. Ver el ejemplo dorado:
%    cursos/1930/1930 Hidráulica Aplicada/Pruebines/pruebin-01-lge-lgp.tex
% \newpage
% ...

% -----------------------------------------------------------------------
% Clave de respuestas -- solo si \toggletrue{icvclave} está activo.
% -----------------------------------------------------------------------
\iftoggle{icvclave}{%
\newpage
\begin{icvnote}[Clave de respuestas]
Desarrollo y resultados numéricos de cada inciso.
\end{icvnote}
}{}

\end{document}
